\documentclass[12pt,a4paper]{article}

% ─── Encoding & fonts ────────────────────────────────────────────────────────
\usepackage[utf8]{inputenc}
\usepackage[T1]{fontenc}
\usepackage{lmodern}

% ─── Page layout ─────────────────────────────────────────────────────────────
\usepackage[margin=2.5cm,top=3cm,bottom=3cm]{geometry}

% ─── Mathematics ─────────────────────────────────────────────────────────────
\usepackage{amsmath}
\usepackage{amssymb}
\usepackage{amsfonts}

% ─── Code listings ───────────────────────────────────────────────────────────
\usepackage{listings}
\usepackage{xcolor}

% ─── Verbatim with LaTeX escapes (used for pseudocode) ───────────────────────
\usepackage{alltt}

% ─── Tables ──────────────────────────────────────────────────────────────────
\usepackage{booktabs}
\usepackage{longtable}
\usepackage{array}
\usepackage{tabularx}

% ─── Miscellaneous ───────────────────────────────────────────────────────────
\usepackage{parskip}
\usepackage{microtype}
\usepackage[hidelinks,colorlinks=true,linkcolor=blue,urlcolor=blue,
            citecolor=blue]{hyperref}
\usepackage{url}
\usepackage{enumitem}

% ─── Colour palette ──────────────────────────────────────────────────────────
\definecolor{codebg}{rgb}{0.97,0.97,0.97}
\definecolor{codeframe}{rgb}{0.80,0.80,0.80}
\definecolor{codegreen}{rgb}{0.00,0.50,0.00}
\definecolor{codeblue}{rgb}{0.00,0.00,0.80}
\definecolor{codepurple}{rgb}{0.50,0.00,0.50}
\definecolor{codegray}{rgb}{0.50,0.50,0.50}

% ─── Listing base style ──────────────────────────────────────────────────────
\lstdefinestyle{generic}{
    backgroundcolor=\color{codebg},
    basicstyle=\small\ttfamily,
    breaklines=true,
    breakatwhitespace=false,
    columns=flexible,
    keepspaces=true,
    showspaces=false,
    showstringspaces=false,
    frame=single,
    rulecolor=\color{codeframe},
    tabsize=4,
}

\lstdefinestyle{python}{
    style=generic,
    language=Python,
    commentstyle=\color{codegreen},
    keywordstyle=\bfseries\color{codeblue},
    stringstyle=\color{codepurple},
    morekeywords={True,False,None,with,as,lambda,yield},
}

\lstdefinestyle{sql}{
    style=generic,
    language=SQL,
    commentstyle=\color{codegreen},
    keywordstyle=\bfseries\color{codeblue},
}

\lstdefinestyle{bash}{
    style=generic,
    language=bash,
    commentstyle=\color{codegreen},
    keywordstyle=\bfseries\color{codeblue},
}

\lstdefinestyle{mermaid}{
    style=generic,
    commentstyle=\color{codegray},
}

\lstset{style=generic}

% ─── Convenience macros ──────────────────────────────────────────────────────
\newcommand{\cs}[1]{\texttt{#1}}          % inline code / monospace
\newcommand{\lrarrow}{$\leftarrow$}        % assignment arrow in prose

% ─────────────────────────────────────────────────────────────────────────────
%  DOCUMENT
% ─────────────────────────────────────────────────────────────────────────────
\begin{document}

% ── Title page ───────────────────────────────────────────────────────────────
\begin{titlepage}
  \centering
  \vspace*{3cm}
  {\huge\bfseries Compressive Sensing Framework\par}
  \vspace{0.8cm}
  {\Large Documentation\par}
  \vspace{2.5cm}
  \begin{tabular}{ll}
    \textbf{Vladislav Gerda}       & \url{https://github.com/hitfot}    \\
    \textbf{Grigory Demchenko}     & \url{https://github.com/Pumukun}   \\
    \textbf{Anastasia Kocherygina} & \url{https://github.com/somniiium} \\
    \textbf{Alina Sabitova}        & \url{https://github.com/AlinaSAB}  \\
    \textbf{Mikhail Shibanov}      & \url{https://github.com/Kar1ch}    \\
  \end{tabular}
  \vfill
  {\normalsize\itshape Compressive Sensing Framework --- Version 1.0\par}
\end{titlepage}

\tableofcontents
\newpage

% =============================================================================
\section{Project Overview}
% =============================================================================

\textbf{Compressive Sensing Framework} is a Python framework for compression
and reconstruction of 2D images using Compressive Sensing (CS) methods.  The
framework implements several sparse signal recovery algorithms along with
quality assessment tools, noise generation utilities, and image smoothing
filters.

\subsection{Key Features}

\begin{itemize}
  \item CS algorithm implementations: \textbf{OMP}, \textbf{CoSaMP},
        \textbf{SP}, \textbf{BRGP}
  \item Discrete Cosine Transform (DCT) as the sparsifying basis matrix
  \item Quality metrics: \textbf{CR} (Compression Ratio), \textbf{PSNR},
        \textbf{SSIM}
  \item Noise generation: Gaussian, Poisson, Salt-and-Pepper, Speckle
  \item Smoothing filters: Mean, Median, Gaussian, Bilateral
  \item Test result persistence in an SQLite database
  \item Result visualisation via Matplotlib
\end{itemize}

% =============================================================================
\section{Theoretical Background}
% =============================================================================

Compressive Sensing allows a sparse signal to be recovered from far fewer
measurements than the Nyquist--Shannon sampling theorem would require.

\subsection{Mathematical Model}

Let $x \in \mathbb{R}^N$ be a sparse signal (an image represented in basis
$\Psi$).  The measurement vector is:

\begin{equation}
  y = \Phi \cdot \Psi \cdot s = \Theta \cdot s
\end{equation}

where:
\begin{itemize}
  \item $\Phi \in \mathbb{R}^{M \times N}$ --- measurement matrix
        (M is much smaller than N)
  \item $\Psi \in \mathbb{R}^{N \times N}$ --- sparsifying basis matrix (DCT)
  \item $s \in \mathbb{R}^N$ --- sparse coefficient vector
  \item $\Theta = \Phi \cdot \Psi$ --- sensing matrix
\end{itemize}

Recovery problem: find $s$ such that $\|s\|_0$ is minimised subject to
$\Theta s \approx y$.

\subsection{Processing Pipeline}

\begin{lstlisting}
+------------------------------------------------------------------+
|                  Compressive Sensing Pipeline                     |
|                                                                    |
|  Image -> Matrix Phi -> y=Phi*x -> CS Algorithm -> xhat -> Psi*xhat |
|  (NxN)      (MxN)        (MxN)      recovery         (NxN)          |
+------------------------------------------------------------------+
\end{lstlisting}

% =============================================================================
\section{Framework Architecture}
% =============================================================================

\subsection{Module Structure}

\begin{lstlisting}
compressive-sensing/
+-- framework/                  # Main framework package
|   +-- __init__.py             # Public API exports
|   +-- omp.py                  # OMP algorithm
|   +-- cosamp.py               # CoSaMP algorithm
|   +-- sp.py                   # SP (Subspace Pursuit) algorithm
|   +-- brgp.py                 # BRGP algorithm
|   +-- transform.py            # DCT transform
|   +-- metrics.py              # Quality metrics (CR, PSNR, SSIM)
|   +-- noise.py                # Noise generation functions
|   +-- smooth.py               # Smoothing filters
|   \-- utils.py                # ImageCS helper class
+-- test/                       # Test scripts
|   +-- omp_test.py             # OMP test
|   +-- cosamp_test.py          # CoSaMP test
|   +-- sp_test.py              # SP test
|   +-- brgp_test.py            # BRGP test
|   +-- all_algs_test.py        # Combined test of all algorithms
|   +-- plot_test.py            # Result visualisation
|   +-- db/                     # Database module
|   |   +-- __init__.py
|   |   +-- db.py               # SQLite CRUD operations
|   |   \-- create_db.sql       # SQL schema
|   \-- plot/                   # Plotting module
|       +-- __init__.py
|       \-- plot.py             # Chart generation
+-- misc/                       # Sample test images
+-- requirements.txt            # Python dependencies
+-- setup.sh                    # Linux setup script
\-- README.md                   # Brief project description
\end{lstlisting}

% ─────────────────────────────────────────────────────────────────────────────
\subsection{Module Dependency Diagram}

The diagram below is written in Mermaid syntax; render it with any
Mermaid-compatible viewer (e.g.\ \url{https://mermaid.live}).

\begin{lstlisting}[style=mermaid,caption={Module dependency graph (Mermaid source)}]
graph TD
    subgraph framework["framework"]
        INIT["__init__.py (public API)"]
        OMP["omp.py"]
        COSAMP["cosamp.py"]
        SP["sp.py"]
        BRGP["brgp.py"]
        DCT["transform.py (DCT)"]
        METRICS["metrics.py (CR / PSNR / SSIM)"]
        NOISE["noise.py"]
        SMOOTH["smooth.py"]
        UTILS["utils.py (ImageCS)"]
    end

    subgraph test["test"]
        OMP_T["omp_test.py"]
        COSAMP_T["cosamp_test.py"]
        SP_T["sp_test.py"]
        BRGP_T["brgp_test.py"]
        ALL_T["all_algs_test.py"]
        PLOT_T["plot_test.py"]
        DB["db/db.py (SQLite)"]
        PLOT["plot/plot.py"]
    end

    INIT --> OMP
    INIT --> COSAMP
    INIT --> SP
    INIT --> BRGP
    INIT --> DCT
    INIT --> UTILS

    OMP --> METRICS
    OMP --> UTILS
    COSAMP --> METRICS
    COSAMP --> UTILS
    SP --> METRICS
    SP --> UTILS
    BRGP --> METRICS
    BRGP --> UTILS
    BRGP --> OMP
    BRGP --> SP

    OMP_T --> INIT
    OMP_T --> DB
    COSAMP_T --> INIT
    COSAMP_T --> DB
    SP_T --> INIT
    SP_T --> DB
    BRGP_T --> INIT
    BRGP_T --> DB
    ALL_T --> INIT
    ALL_T --> DB
    PLOT_T --> PLOT
    PLOT --> DB
\end{lstlisting}

% ─────────────────────────────────────────────────────────────────────────────
\subsection{Class Diagram}

\begin{lstlisting}[style=mermaid,caption={UML class diagram (Mermaid source)}]
classDiagram
    class ImageCS {
        -np.ndarray __matrix
        -float __cr
        -float __psnr
        -float __ssim
        +__init__(matrix, cr, psnr, ssim)
        +get_Image() np.ndarray
        +get_CR() float
        +get_PSNR() float
        +get_SSIM() float
        +set_Image(image) None
        +set_CR(cr) None
        +set_PSNR(psnr) None
        +set_SSIM(ssim) None
    }

    class OMP {
        +omp(image_path, matrix, M, K) ImageCS
        +cs_omp(y, Phi, K) Tuple
    }

    class CoSaMP {
        +cosamp(image_path, matrix, s, M) ImageCS
        +cs_cosamp(y, s, Phi, epsilon, K) np.ndarray
    }

    class SP {
        +sp(image_path, matrix, M, K) ImageCS
        +cs_sp(y, Phi, K) Tuple
    }

    class BRGP {
        +brgp(image_path, matrix, M, K) ImageCS
        +cs_brgp(y, Phi, K, Candidate, u) np.ndarray
    }

    class Transform {
        +dct(N) np.ndarray
    }

    class Metrics {
        +PSNR(original, compressed) float
        +SSIM(original, compressed) float
        +CR(image_source, image_compressed) float
    }

    class Noise {
        +GaussianNoise(image_path, stdev, show) np.ndarray
        +PoissonNoise(image_path, show) np.ndarray
        +SaltAndPepperNoise(image_path, show, number_of_pixels) np.ndarray
        +SpeckleNoise(image_path, show, variance) np.ndarray
    }

    class Smooth {
        +Mean_filter(image_path, k, show) np.ndarray
        +Median_filter(image_path, k, show) np.ndarray
        +Gaussian_filter(image_path, k, show) np.ndarray
        +Bilateral_filter(image_path, k, show) np.ndarray
    }

    class Database {
        +connect_db() Connection
        +create_table() None
        +add_result(...) None
        +get_all_results() List
        +get_result_by_id(id) Tuple
        +get_result_by_alg(alg) List
        +update_result(...) None
        +delete_result(alg) None
        +delete_all() None
    }

    OMP --> ImageCS : returns
    CoSaMP --> ImageCS : returns
    SP --> ImageCS : returns
    BRGP --> ImageCS : returns
    BRGP --> OMP : uses cs_omp
    BRGP --> SP : uses cs_sp
    OMP --> Metrics : computes CR/PSNR
    CoSaMP --> Metrics : computes CR/PSNR/SSIM
    SP --> Metrics : computes CR/PSNR
    BRGP --> Metrics : computes CR/PSNR
\end{lstlisting}

% ─────────────────────────────────────────────────────────────────────────────
\subsection{Data Flow}

\begin{lstlisting}[style=mermaid,caption={End-to-end data flow (Mermaid source)}]
flowchart LR
    A["Input image / image_path"] --> B["Load as grayscale (cv2.imread)"]
    B --> C["Image matrix: NxN numpy array"]
    C --> D["Measurement matrix Phi ~ N(0,1/M) : MxN"]
    C --> E["Basis matrix Psi = DCT : NxN"]
    D & E --> F["Sensing matrix Theta = Phi*Psi : MxN"]
    C & D --> G["Measurement vector y = Phi*im : MxW"]
    G & F --> H{"CS Algorithm (column by column)"}
    H --> |OMP|    I1["cs_omp(y, Theta, K)"]
    H --> |CoSaMP| I2["cs_cosamp(y, s, Theta)"]
    H --> |SP|     I3["cs_sp(y, Theta, K)"]
    H --> |BRGP|   I4["cs_brgp(y, Theta, K, Candidate)"]
    I1 & I2 & I3 & I4 --> J["Sparse matrix sparse_rec_1d: NxN"]
    J --> K["Reconstructed image: img_rec = Psi * sparse_rec_1d"]
    K --> L["Quality metrics: CR / PSNR / SSIM"]
    K & L --> M["ImageCS: get_Image / get_CR / get_PSNR / get_SSIM"]
\end{lstlisting}

% =============================================================================
\section{Framework Components}
% =============================================================================

\subsection{Recovery Algorithms}

All algorithms work with 2D images using a \textbf{column-by-column} approach:
each pixel column is recovered independently as a 1D CS problem.

% ─────────────────────────────────────────────────────────────────────────────
\subsubsection{OMP --- Orthogonal Matching Pursuit}

\textbf{File:} \cs{framework/omp.py}\quad\textbf{Author:} Vladislav Gerda

\textbf{Description.}  OMP is a greedy algorithm that iteratively selects the
column of the sensing matrix most correlated with the current residual, then
projects the measured signal onto the selected subspace.

\textbf{Pseudocode:}

\begin{alltt}
\textbf{function} cs\_omp(y, \(\Phi\), K):
    residual \(\leftarrow\) y
    index    \(\leftarrow\) array of \(-1\), length N        \textit{// support mask}

    \textbf{for} j = 1 \textbf{to} K:
        product    \(\leftarrow\) |\(\Phi^T \cdot\) residual|  \textit{// correlation with each column}
        pos        \(\leftarrow\) argmax(product)      \textit{// most correlated column}
        index[pos] \(\leftarrow\) 1                    \textit{// add to support}

        a        \(\leftarrow\) pinv(\(\Phi\)[:, index\(\geq\)0]) \(\cdot\) y  \textit{// least-squares on support}
        residual \(\leftarrow\) y \(-\) \(\Phi\)[:, index\(\geq\)0] \(\cdot\) a   \textit{// update residual}

    result[index\(\geq\)0] \(\leftarrow\) a
    \textbf{return} result
\end{alltt}

\textbf{Parameters:}

\begin{center}
\begin{tabular}{lll}
  \toprule
  Parameter & Type & Description \\
  \midrule
  \cs{image\_path} & \cs{str}          & Path to the input image \\
  \cs{matrix}      & \cs{np.ndarray}   & Basis matrix $N \times N$ (DCT) \\
  \cs{M}           & \cs{int}          & Number of measurements ($M < N$) \\
  \cs{K}           & \cs{int}          & Number of iterations (= sparsity level) \\
  \bottomrule
\end{tabular}
\end{center}

\textbf{Complexity:} $\mathcal{O}(K \cdot M \cdot N)$ per column.

\bigskip\hrule\bigskip

% ─────────────────────────────────────────────────────────────────────────────
\subsubsection{CoSaMP --- Compressive Sampling Matching Pursuit}

\textbf{File:} \cs{framework/cosamp.py}\quad\textbf{Author:} Vladislav Gerda

\textbf{Description.}  CoSaMP extends OMP by selecting $2s$ best candidate
indices at each step, merging them with the current support, solving a
least-squares problem, and pruning back to the $s$ largest components.  The
algorithm iterates until convergence.

\textbf{Pseudocode:}

\begin{alltt}
\textbf{function} cs\_cosamp(y, s, \(\Phi\), \(\varepsilon\)=1e-10, K=1000):
    residual \(\leftarrow\) y
    result   \(\leftarrow\) zeros(N)

    \textbf{for} j = 1 \textbf{to} K:
        product   \(\leftarrow\) |\(\Phi^T \cdot\) residual|
        top\_k\_idx \(\leftarrow\) indices of 2s largest values in product
        top\_k\_idx \(\leftarrow\) top\_k\_idx \(\cup\) nonzero(result)  \textit{// merge with current support}

        x           \(\leftarrow\) zeros(N)
        x[top\_k\_idx] \(\leftarrow\) lstsq(\(\Phi\)[:, top\_k\_idx], y)  \textit{// constrained least squares}
        set all but the s largest-magnitude entries of x to 0

        result       \(\leftarrow\) x
        residual\_old \(\leftarrow\) residual
        residual     \(\leftarrow\) y \(-\) \(\Phi \cdot\) result

        \textbf{if} \(\|\)residual\(\|\) < \(\varepsilon\)  \textbf{or}  \(\|\)residual \(-\) residual\_old\(\|\) < \(\varepsilon\):
            \textbf{break}

    \textbf{return} |result|
\end{alltt}

\textbf{Parameters:}

\begin{center}
\begin{tabular}{lll}
  \toprule
  Parameter   & Type            & Description \\
  \midrule
  \cs{image\_path} & \cs{str}       & Path to the input image \\
  \cs{matrix}      & \cs{np.ndarray}& Basis matrix $N \times N$ (DCT) \\
  \cs{s}           & \cs{int}       & Signal sparsity level \\
  \cs{M}           & \cs{int}       & Number of measurements \\
  \cs{epsilon}     & \cs{float}     & Convergence tolerance (default \cs{1e-10}) \\
  \cs{K}           & \cs{int}       & Maximum number of iterations (default \cs{1000}) \\
  \bottomrule
\end{tabular}
\end{center}

\bigskip\hrule\bigskip

% ─────────────────────────────────────────────────────────────────────────────
\subsubsection{SP --- Subspace Pursuit}

\textbf{File:} \cs{framework/sp.py}\quad\textbf{Author:} Grigory Demchenko

\textbf{Description.}  SP maintains a fixed support set of size $K$ at every
iteration: it expands the set with $K$ new candidates, solves a constrained LS
problem over the union, then retains only the $K$ largest-magnitude components.

\textbf{Pseudocode:}

\begin{alltt}
\textbf{function} cs\_sp(y, \(\Phi\), K):
    residual \(\leftarrow\) y
    index    \(\leftarrow\) \(\emptyset\)                       \textit{// current support set}
    x        \(\leftarrow\) zeros(N)

    \textbf{for} j = 1 \textbf{to} K:
        product   \(\leftarrow\) |\(\Phi^T \cdot\) residual|
        top\_k\_idx \(\leftarrow\) indices of K largest values in product
        index     \(\leftarrow\) index \(\cup\) top\_k\_idx      \textit{// expand support}

        x\_temp    \(\leftarrow\) pinv(\(\Phi\)[:, index]) \(\cdot\) y  \textit{// least squares on support}
        x[index]  \(\leftarrow\) x\_temp

        index    \(\leftarrow\) indices of K largest values in |x|  \textit{// prune to K}
        residual \(\leftarrow\) y \(-\) \(\Phi \cdot\) x

    \textbf{return} x, index
\end{alltt}

\textbf{Parameters:}

\begin{center}
\begin{tabular}{lll}
  \toprule
  Parameter        & Type             & Description \\
  \midrule
  \cs{image\_path} & \cs{str}         & Path to the input image \\
  \cs{matrix}      & \cs{np.ndarray}  & Basis matrix $N \times N$ (DCT) \\
  \cs{M}           & \cs{int}         & Number of measurements \\
  \cs{K}           & \cs{int}         & Number of iterations and support set size \\
  \bottomrule
\end{tabular}
\end{center}

\bigskip\hrule\bigskip

% ─────────────────────────────────────────────────────────────────────────────
\subsubsection{BRGP --- Backtracking Refined Greedy Pursuit}

\textbf{File:} \cs{framework/brgp.py}\quad\textbf{Author:} Grigory Demchenko

\textbf{Description.}  BRGP is a hybrid algorithm.  It initialises the support
as the intersection of the SP and OMP candidate sets, then iteratively refines
the support with a backtracking mechanism (reverting to the last saved state
when the residual worsens).  A final SP-style refinement phase shrinks the
support from size $K$ down to 0.

\textbf{Pseudocode:}

\begin{alltt}
\textbf{function} brgp(y, \(\Phi\), K, u=0.8):
    \textit{// --- Initialisation ---}
    \_, Candidate\_sp  \(\leftarrow\) cs\_sp(y, \(\Phi\), K)
    \_, Candidate\_omp \(\leftarrow\) cs\_omp(y, \(\Phi\), K)
    Candidate         \(\leftarrow\) Candidate\_sp \(\cap\) Candidate\_omp  \textit{// intersection as seed}

    x      \(\leftarrow\) pinv(\(\Phi\)[:, Candidate]) \(\cdot\) y projected onto Candidate
    r      \(\leftarrow\) y \(-\) \(\Phi \cdot\) x
    r\_save \(\leftarrow\) r
    C\_save \(\leftarrow\) Candidate

    \textit{// --- Expansion phase ---}
    F         \(\leftarrow\) \(\{\) i : |\(\Phi^T r\)|\_i > u \(\cdot\) max|\(\Phi^T r\)| \(\}\)
    Candidate \(\leftarrow\) Candidate \(\cup\) F
    x, r      \(\leftarrow\) recompute(\(\Phi\), Candidate, y)

    \textbf{while} len(Candidate) < K:
        \textbf{if} \(\|\)r \(-\) r\_save\(\|\) < \(\|\)y\(\|\):           \textit{// improvement: expand}
            C\_save    \(\leftarrow\) Candidate
            r\_save    \(\leftarrow\) r
            F         \(\leftarrow\) \(\{\) i : |\(\Phi^T r\)|\_i > u \(\cdot\) max|\(\Phi^T r\)| \(\}\)
            Candidate \(\leftarrow\) Candidate \(\cup\) F
        \textbf{else}:                              \textit{// no improvement: backtrack}
            Candidate \(\leftarrow\) C\_save
            r         \(\leftarrow\) r\_save
            C\_dif     \(\leftarrow\) complement candidates not yet in Candidate
            F         \(\leftarrow\) argmax of pinv(\(\Phi\)[:, C\_dif]) \(\cdot\) y
            Candidate \(\leftarrow\) Candidate \(\cup\) F
        x, r \(\leftarrow\) recompute(\(\Phi\), Candidate, y)

    \textit{// --- SP-style refinement phase ---}
    T \(\leftarrow\) K
    \textbf{while} T > 0:
        top\_T     \(\leftarrow\) indices of T largest values in |\(\Phi^T r\)|
        Candidate \(\leftarrow\) Candidate \(\cup\) top\_T
        keep\_K    \(\leftarrow\) indices of K largest values in |pinv(\(\Phi\)[:, Candidate]) \(\cdot\) y|
        Candidate \(\leftarrow\) Candidate[keep\_K]
        x, r      \(\leftarrow\) recompute(\(\Phi\), Candidate, y)
        T         \(\leftarrow\) floor(T \(\cdot\) u)

    \textbf{return} x
\end{alltt}

\textbf{Parameters:}

\begin{center}
\begin{tabular}{lll}
  \toprule
  Parameter        & Type            & Description \\
  \midrule
  \cs{image\_path} & \cs{str}        & Path to the input image \\
  \cs{matrix}      & \cs{np.ndarray} & Basis matrix $N \times N$ (DCT) \\
  \cs{M}           & \cs{int}        & Number of measurements \\
  \cs{K}           & \cs{int}        & Number of iterations \\
  \cs{u}           & \cs{float}      & Backtracking coefficient ($0 < u < 1$, default 0.8) \\
  \bottomrule
\end{tabular}
\end{center}

\bigskip\hrule\bigskip

% ─────────────────────────────────────────────────────────────────────────────
\subsection{Algorithm Comparison}

\begin{center}
\begin{tabular}{lllll}
  \toprule
  Property                & OMP      & CoSaMP            & SP       & BRGP     \\
  \midrule
  Type                    & Greedy   & Iterative         & Iterative & Hybrid   \\
  Convergence guarantee   & No       & Yes               & Yes       & Partial  \\
  Iteration count         & Fixed $K$& Until convergence & Fixed $K$ & Adaptive \\
  Uses other algorithms   & No       & No                & No        & OMP + SP \\
  Available metrics       & CR, PSNR & CR, PSNR, SSIM    & CR, PSNR  & CR, PSNR \\
  Computational cost      & Medium   & High              & Medium    & High     \\
  \bottomrule
\end{tabular}
\end{center}

% ─────────────────────────────────────────────────────────────────────────────
\subsection{Transforms}

\subsubsection{DCT --- Discrete Cosine Transform}

\textbf{File:} \cs{framework/transform.py}

\begin{lstlisting}[style=python]
def dct(N: int) -> np.ndarray
\end{lstlisting}

Generates an orthonormal DCT matrix of size $N \times N$.  Used as the
sparsifying basis matrix $\Psi$ --- natural images typically have very few
non-zero DCT coefficients.

\textbf{Formula for column $k$:}

\[
  \psi_k[n] = \cos\!\left(n \cdot \frac{k\pi}{N}\right), \quad n = 0, \ldots, N-1
\]
\[
  \psi_k = \frac{\psi_k - \overline{\psi_k}}{\|\psi_k\|}, \quad k > 0
\]

\textbf{Example:}

\begin{lstlisting}[style=python]
from framework import dct
Psi = dct(256)  # 256x256 basis matrix
\end{lstlisting}

% ─────────────────────────────────────────────────────────────────────────────
\subsection{Quality Metrics}

\textbf{File:} \cs{framework/metrics.py}

\subsubsection{CR --- Compression Ratio}

\begin{lstlisting}[style=python]
def CR(image_source: np.ndarray, image_compressed: np.ndarray) -> float
\end{lstlisting}

Ratio of the number of non-zero elements in the original image to the number
of non-zero elements in its sparse representation.

\[
  \mathrm{CR} = \frac{\#\,\mathrm{nonzero}(x)}{\#\,\mathrm{nonzero}(\hat{x})}
\]

\begin{quote}
  \textbf{Note:} CR $> 1$ means the sparse representation has fewer non-zero
  elements, i.e.\ the signal is genuinely sparse.
\end{quote}

\subsubsection{PSNR --- Peak Signal-to-Noise Ratio}

\begin{lstlisting}[style=python]
def PSNR(original: np.ndarray, compressed: np.ndarray) -> float
\end{lstlisting}

Peak signal-to-noise ratio in dB.  Uses
\cs{skimage.metrics.peak\_signal\_noise\_ratio}.  Higher values indicate better
reconstruction quality; values $> 30\,\mathrm{dB}$ are generally considered
acceptable.

\subsubsection{SSIM --- Structural Similarity Index}

\begin{lstlisting}[style=python]
def SSIM(original: np.ndarray, compressed: np.ndarray) -> float
\end{lstlisting}

Structural Similarity Index (0--1).  Uses
\cs{skimage.metrics.structural\_similarity}.  Values close to 1 indicate high
structural fidelity to the original.

% ─────────────────────────────────────────────────────────────────────────────
\subsection{Noise Functions}

\textbf{File:} \cs{framework/noise.py}

\begin{center}
\begin{tabular}{p{6cm}p{4cm}p{4.5cm}}
  \toprule
  Function & Description & Parameters \\
  \midrule
  \cs{GaussianNoise(image\_path, stdev, show)}
    & Gaussian blur (\cs{cv2.GaussianBlur})
    & \cs{stdev} --- kernel size (odd integer: 3, 5, 7, \ldots) \\[4pt]
  \cs{PoissonNoise(image\_path, show)}
    & Poisson noise
    & --- \\[4pt]
  \cs{SaltAndPepperNoise(image\_path, show, number\_of\_pixels)}
    & Salt-and-pepper noise
    & \cs{number\_of\_pixels} --- number of corrupted pixels \\[4pt]
  \cs{SpeckleNoise(image\_path, show, variance)}
    & Speckle (multiplicative) noise
    & \cs{variance} --- noise variance \\
  \bottomrule
\end{tabular}
\end{center}

All functions accept \verb|show: bool = True| to display a side-by-side
comparison of the original and noisy image via Matplotlib.

% ─────────────────────────────────────────────────────────────────────────────
\subsection{Smoothing Filters}

\textbf{File:} \cs{framework/smooth.py}

\begin{center}
\begin{tabular}{p{5.5cm}p{4cm}p{4.5cm}}
  \toprule
  Function & Description & Parameters \\
  \midrule
  \cs{Mean\_filter(image\_path, k, show)}
    & Averaging filter
    & \cs{k} --- kernel size ($k \times k$) \\[4pt]
  \cs{Median\_filter(image\_path, k, show)}
    & Median filter
    & \cs{k} --- kernel size \\[4pt]
  \cs{Gaussian\_filter(image\_path, k, show)}
    & Gaussian filter
    & \cs{k} --- kernel size (must be odd) \\[4pt]
  \cs{Bilateral\_filter(image\_path, k, show)}
    & Bilateral filter
    & \cs{k} --- pixel neighbourhood diameter \\
  \bottomrule
\end{tabular}
\end{center}

% ─────────────────────────────────────────────────────────────────────────────
\subsection{ImageCS Class}

\textbf{File:} \cs{framework/utils.py}

A container that holds the result of a CS algorithm: the reconstructed image
together with quality metrics.

\begin{lstlisting}[style=python]
class ImageCS:
    def __init__(self, matrix: np.ndarray, cr: float = 0.0,
                 psnr: float = 0.0, ssim: float = 0.0)

    def get_Image(self) -> np.ndarray   # Image data
    def get_CR(self) -> float           # Compression ratio
    def get_PSNR(self) -> float         # PSNR in dB
    def get_SSIM(self) -> float         # SSIM (0-1)

    def set_Image(self, image: np.ndarray) -> None
    def set_CR(self, cr: float) -> None
    def set_PSNR(self, psnr: float) -> None
    def set_SSIM(self, ssim: float) -> None
\end{lstlisting}

% =============================================================================
\section{Testing Infrastructure}
% =============================================================================

\subsection{Test Structure}

Test scripts are located in the \cs{test/} directory and must be
\textbf{run from within that directory}.

\begin{lstlisting}
test/
+-- omp_test.py       # OMP test: sweeps M and K, saves images and metrics
+-- cosamp_test.py    # CoSaMP test
+-- sp_test.py        # SP test
+-- brgp_test.py      # BRGP test
+-- all_algs_test.py  # Parallel test of all algorithms (threading)
+-- plot_test.py      # Plot results from DB
+-- db/               # SQLite module
\-- plot/             # Plotting module
\end{lstlisting}

\subsection{Database (SQLite)}

\textbf{File:} \cs{test/db/db.py}

Schema of the \cs{results} table:

\begin{lstlisting}[style=sql]
CREATE TABLE IF NOT EXISTS results (
    id             INTEGER PRIMARY KEY AUTOINCREMENT,
    original_image TEXT    NOT NULL,   -- source image name
    pwd            TEXT    NOT NULL,   -- path to reconstructed image
    algorithm      TEXT    NOT NULL,   -- algorithm name
    PSNR           FLOAT,              -- PSNR metric
    SSIM           FLOAT,              -- SSIM metric
    CR             FLOAT,              -- compression ratio
    K              INTEGER NOT NULL,   -- iteration count
    M              INTEGER NOT NULL,   -- measurement matrix size
    height         INTEGER NOT NULL,   -- image height
    width          INTEGER NOT NULL    -- image width
);
\end{lstlisting}

\subsection{Running Tests}

\begin{lstlisting}[style=bash]
cd test/

# Test a single algorithm
python omp_test.py

# Test all algorithms in parallel
python all_algs_test.py

# Plot results from the database
python plot_test.py
\end{lstlisting}

% =============================================================================
\section{Integration Guide}
% =============================================================================

\subsection{Installation}

\textbf{Linux:}
\begin{lstlisting}[style=bash]
git clone <repo_url>
cd compressive-sensing
./setup.sh
\end{lstlisting}

\textbf{Windows:}
\begin{lstlisting}
python -m venv venv
venv\Scripts\activate
pip install -r requirements.txt
\end{lstlisting}

\subsection{Basic Usage}

\begin{lstlisting}[style=python]
import sys
sys.path.insert(0, '/path/to/compressive-sensing')

from framework import omp, cosamp, sp, brgp, dct

# 1. Build a DCT basis matrix matching the image size
#    (N = image height / width)
N = 256
Psi = dct(N)

# 2. Choose algorithm parameters
M = 128   # Number of measurements (M < N; M ~ N/2 recommended)
K = 20    # Iterations / sparsity level

# 3. Run the algorithm
result = omp("misc/lena.png", Psi, M, K)

# 4. Retrieve results
import cv2
cv2.imwrite("output.png", result.get_Image())
print(f"CR:   {result.get_CR():.3f}")
print(f"PSNR: {result.get_PSNR():.2f} dB")
\end{lstlisting}

\subsection{Example: Multiple Algorithms}

\begin{lstlisting}[style=python]
from framework import omp, cosamp, sp, brgp, dct
import cv2

image_path = "misc/lena.png"
N = 256
Psi = dct(N)
M = 128
K = 20

algorithms = {
    "OMP":    lambda: omp(image_path, Psi, M, K),
    "SP":     lambda: sp(image_path, Psi, M, K),
    "BRGP":   lambda: brgp(image_path, Psi, M, K),
    "CoSaMP": lambda: cosamp(image_path, Psi, K, M),  # s=K for CoSaMP
}

for name, alg_fn in algorithms.items():
    result = alg_fn()
    cv2.imwrite(f"output_{name.lower()}.png", result.get_Image())
    print(f"{name}: CR={result.get_CR():.3f}, PSNR={result.get_PSNR():.2f} dB")
\end{lstlisting}

\subsection{Example: Noise and Filtering Pre-processing}

\begin{lstlisting}[style=python]
import cv2
import numpy as np
from framework.noise import GaussianNoise
from framework.smooth import Median_filter
from framework import omp, dct

# Add noise
noisy = GaussianNoise("misc/lena.png", stdev=5, show=False)
cv2.imwrite("/tmp/noisy.png", noisy)

# Apply filter
filtered = Median_filter("/tmp/noisy.png", k=3, show=False)
cv2.imwrite("/tmp/filtered.png", filtered)

# Apply CS to the filtered image
result = omp("/tmp/filtered.png", dct(256), M=128, K=20)
print(f"PSNR after noise + filter + CS: {result.get_PSNR():.2f} dB")
\end{lstlisting}

\subsection{Example: Using Metrics Directly}

\begin{lstlisting}[style=python]
import cv2
import numpy as np
import framework.metrics as metrics

original   = cv2.imread("misc/lena.png", cv2.IMREAD_GRAYSCALE)
compressed = cv2.imread("output.png",    cv2.IMREAD_GRAYSCALE)

psnr_val = metrics.PSNR(original, compressed)
ssim_val = metrics.SSIM(original, compressed)
cr_val   = metrics.CR(original, compressed)

print(f"PSNR: {psnr_val:.2f} dB")
print(f"SSIM: {ssim_val:.4f}")
print(f"CR:   {cr_val:.3f}")
\end{lstlisting}

\subsection{Example: Saving Results to the Database}

Run from the \cs{test/} directory:

\begin{lstlisting}[style=python]
import sys
sys.path.insert(0, '/path/to/compressive-sensing/test')
import db

db.create_table()
db.add_result(
    pwd="output_omp.png",
    original_image="lena",
    algorithm="OMP",
    psnr=32.5,
    ssim=0.91,
    cr=1.75,
    k=20,
    m=128,
    height=256,
    width=256
)
results = db.get_all_results()
\end{lstlisting}

% =============================================================================
\section{Scaling}
% =============================================================================

\subsection{Current Limitations}

\begin{center}
\begin{tabular}{p{5cm}p{9cm}}
  \toprule
  Limitation & Description \\
  \midrule
  Images assumed square
    & All algorithms set $N = H$ (image height); pass \cs{dct(H)} for
      rectangular images. \\[4pt]
  Grayscale only
    & All algorithms load images with \cs{IMREAD\_GRAYSCALE}. \\[4pt]
  Sequential column processing
    & Each column is processed in a \cs{for i in range(W)} loop. \\[4pt]
  No GPU acceleration
    & All computation runs on CPU via NumPy. \\
  \bottomrule
\end{tabular}
\end{center}

\begin{quote}
  \textbf{Rectangular images:} $\Phi$ is built over the height dimension
  ($N = H$) and the loop \verb|for i in range(W)| covers all columns
  regardless of width, so rectangular images work correctly.  Always pass
  \cs{dct(H)} rather than \cs{dct(W)}.
\end{quote}

\subsection{Horizontal Scaling}

\subsubsection{Column-level Parallelism}

The current sequential column loop can easily be parallelised:

\begin{lstlisting}[style=python]
from concurrent.futures import ThreadPoolExecutor
import numpy as np

def process_column(i, y_col, Theta, K):
    y = np.reshape(y_col, (-1, 1))
    col_rec, _ = cs_omp(y, Theta, K)
    return i, np.reshape(col_rec, (-1,))

with ThreadPoolExecutor(max_workers=8) as executor:
    futures = [
        executor.submit(process_column, i, img_cs_1d[:, i], Theta_1d, K)
        for i in range(W)
    ]
    for future in futures:
        i, col = future.result()
        sparse_rec_1d[:, i] = col
\end{lstlisting}

\subsubsection{Image-level Parallelism}

For batch processing of multiple images use \cs{threading} (as in
\cs{all\_algs\_test.py}) or \cs{multiprocessing}:

\begin{lstlisting}[style=python]
from multiprocessing import Pool
from framework import omp, dct

def process_image(args):
    image_path, M, K = args
    return omp(image_path, dct(256), M, K)

images = ["misc/lena.png", "misc/house.png", "misc/4.1.05.png"]
params = [(img, 128, 20) for img in images]

with Pool(processes=4) as pool:
    results = pool.map(process_image, params)
\end{lstlisting}

\subsubsection{GPU Acceleration via CuPy}

NumPy operations can be offloaded to a GPU with minimal code changes:

\begin{lstlisting}[style=python]
import cupy as cp   # pip install cupy-cuda12x
import numpy as np

# Replace np.dot -> cp.dot, np.linalg -> cp.linalg
Phi_gpu    = cp.array(Phi)
matrix_gpu = cp.array(matrix)
img_gpu    = cp.array(im)

img_cs_1d = cp.dot(Phi_gpu, img_gpu)
Theta_1d  = cp.dot(Phi_gpu, matrix_gpu)
# ... rest of the algorithm on GPU ...
result_np = cp.asnumpy(sparse_rec_1d)
\end{lstlisting}

\subsection{Vertical Scaling}

\subsubsection{Colour Image Support}

Extend to colour (RGB) images by processing each channel independently:

\begin{lstlisting}[style=python]
def omp_color(image_path: str, matrix: np.ndarray, M: int, K: int) -> np.ndarray:
    image = cv2.imread(image_path)  # BGR
    channels = cv2.split(image)
    rec_channels = []
    for ch in channels:
        # Save the channel to a temp file, or pass as ndarray directly
        result = _omp_channel(ch, matrix, M, K)
        rec_channels.append(result)
    return cv2.merge(rec_channels)
\end{lstlisting}

\subsubsection{Rectangular Image Support}

For images where $H \neq W$, apply 2D DCT separately along rows and columns,
or pass two different matrices:

\begin{lstlisting}[style=python]
# Option: run the algorithm on the transposed matrix
# to process rows, then columns.
\end{lstlisting}

\subsubsection{Deterministic Measurement Matrix}

Replace the random Gaussian $\Phi$ with a deterministic matrix (Hadamard,
Toeplitz) for reproducibility:

\begin{lstlisting}[style=python]
from scipy.linalg import hadamard

N   = 256
H   = hadamard(N)
Phi = H[:M, :] / np.sqrt(M)  # First M rows of the Hadamard matrix
\end{lstlisting}

% ─────────────────────────────────────────────────────────────────────────────
\subsection{Adding a New Algorithm}

The framework is designed for easy extension.  To add a new algorithm:

\textbf{Step 1.} Create \cs{framework/my\_algorithm.py}:

\begin{lstlisting}[style=python]
import numpy as np
import cv2
import framework.metrics as metrics
from framework.utils import ImageCS
from typing import Tuple

def my_algorithm(image_path: str, matrix: np.ndarray, M: int, K: int) -> ImageCS:
    """
    Algorithm description.
        image_path - path to the image.
        matrix     - NxN basis matrix.
        M          - number of measurements.
        K          - number of iterations.
    """
    image = cv2.imread(image_path, cv2.IMREAD_GRAYSCALE)
    H, W = image.shape
    N = H

    im = np.array(image)
    Phi = np.random.randn(M, N) / np.sqrt(M)
    img_cs_1d = np.dot(Phi, im)
    Theta_1d = np.dot(Phi, matrix)
    sparse_rec_1d = np.zeros((N, W))

    for i in range(W):
        y = np.reshape(img_cs_1d[:, i], (M, 1))
        column_rec = _cs_my_algorithm(y, Theta_1d, K)
        sparse_rec_1d[:, i] = np.reshape(column_rec, (N,))

    img_rec = np.dot(matrix, sparse_rec_1d)

    CR   = metrics.CR(image, sparse_rec_1d)
    PSNR = metrics.PSNR(image, img_rec)

    return ImageCS(img_rec, cr=CR, psnr=PSNR)


def _cs_my_algorithm(y: np.ndarray, Phi: np.ndarray, K: int) -> np.ndarray:
    """Core algorithm for a single column."""
    # ... your implementation ...
    pass
\end{lstlisting}

\textbf{Step 2.} Register it in \cs{framework/\_\_init\_\_.py}:

\begin{lstlisting}[style=python]
from .my_algorithm import my_algorithm
\end{lstlisting}

\textbf{Step 3.} Add a test script \cs{test/my\_algorithm\_test.py} following
the pattern of the existing test files.

% ─────────────────────────────────────────────────────────────────────────────
\subsection{Parameter Selection Guide}

Use the following tables to choose starting values for $M$ and $K$ based on
image size and the speed/quality trade-off.

\begin{center}
\begin{tabular}{llll}
  \toprule
  Image size $N$ & Recommended $M$ & Recommended $K$ & Priority \\
  \midrule
  $N \leq 256$          & 128 ($N/2$)   & 10--30  & ---      \\
  $256 < N \leq 512$    & 256 ($N/2$)   & 20--50  & ---      \\
  $N > 512$             & $N/3$--$N/2$  & 30--100 & ---      \\
  Any                   & $< N/2$       & 10--20  & Speed    \\
  Any                   & $\approx N$   & 50--200 & Quality  \\
  Any                   & $N/2$         & 20--50  & Balanced \\
  \bottomrule
\end{tabular}
\end{center}

\begin{center}
\begin{tabular}{lll}
  \toprule
  Parameter & Recommended range & Effect \\
  \midrule
  \cs{M} & $N/4$--$N$ & $\uparrow M \Rightarrow \uparrow\mathrm{PSNR},\;\uparrow\mathrm{time}$ \\
  \cs{K} (OMP/SP/BRGP) & 10--200 & $\uparrow K \Rightarrow \uparrow\mathrm{PSNR}$ (up to saturation), $\uparrow\mathrm{time}$ \\
  \cs{s} (CoSaMP) & 5--50 & Signal sparsity \\
  \cs{u} (BRGP) & 0.6--0.9 & Expansion aggressiveness \\
  \bottomrule
\end{tabular}
\end{center}

% =============================================================================
\section{API Reference}
% =============================================================================

\subsection{\texttt{framework} (public API)}

\begin{lstlisting}[style=python]
from framework import omp, cosamp, sp, brgp, dct, ImageCS
\end{lstlisting}

\begin{description}[style=nextline]
  \item[\cs{omp(image\_path, matrix, M, K) $\to$ ImageCS}]
    Image reconstruction using Orthogonal Matching Pursuit.
  \item[\cs{cosamp(image\_path, matrix, s, M) $\to$ ImageCS}]
    Image reconstruction using Compressive Sampling Matching Pursuit.
  \item[\cs{sp(image\_path, matrix, M, K) $\to$ ImageCS}]
    Image reconstruction using Subspace Pursuit.
  \item[\cs{brgp(image\_path, matrix, M, K) $\to$ ImageCS}]
    Image reconstruction using Backtracking Refined Greedy Pursuit.
  \item[\cs{dct(N) $\to$ np.ndarray}]
    Generate an orthonormal DCT matrix of size $N \times N$.
\end{description}

\subsection{\texttt{framework.metrics}}

\begin{lstlisting}[style=python]
import framework.metrics as metrics
\end{lstlisting}

\begin{description}
  \item[\cs{metrics.PSNR(original, compressed) $\to$ float}] ~
  \item[\cs{metrics.SSIM(original, compressed) $\to$ float}] ~
  \item[\cs{metrics.CR(image\_source, image\_compressed) $\to$ float}] ~
\end{description}

\subsection{\texttt{framework.noise}}

\begin{lstlisting}[style=python]
from framework.noise import (GaussianNoise, PoissonNoise,
                              SaltAndPepperNoise, SpeckleNoise)
\end{lstlisting}

\begin{description}
  \item[\cs{GaussianNoise(image\_path, stdev=5, show=True) $\to$ np.ndarray}] ~
  \item[\cs{PoissonNoise(image\_path, show=True) $\to$ np.ndarray}] ~
  \item[\cs{SaltAndPepperNoise(image\_path, show=True, number\_of\_pixels=1000) $\to$ np.ndarray}] ~
  \item[\cs{SpeckleNoise(image\_path, show=True, variance=0.1) $\to$ np.ndarray}] ~
\end{description}

\subsection{\texttt{framework.smooth}}

\begin{lstlisting}[style=python]
from framework.smooth import (Mean_filter, Median_filter,
                               Gaussian_filter, Bilateral_filter)
\end{lstlisting}

\begin{description}
  \item[\cs{Mean\_filter(image\_path, k, show=True) $\to$ np.ndarray}] ~
  \item[\cs{Median\_filter(image\_path, k, show=True) $\to$ np.ndarray}] ~
  \item[\cs{Gaussian\_filter(image\_path, k, show=True) $\to$ np.ndarray}] ~
  \item[\cs{Bilateral\_filter(image\_path, k, show=True) $\to$ np.ndarray}] ~
\end{description}

\subsection{\texttt{test/db}}

\begin{lstlisting}[style=python]
import db  # from the test/ directory
\end{lstlisting}

\begin{description}
  \item[\cs{db.create\_table() $\to$ None}] ~
  \item[\cs{db.add\_result(pwd, original\_image, algorithm, psnr, ssim, cr, k, m, height, width) $\to$ None}] ~
  \item[\cs{db.get\_all\_results() $\to$ List[Tuple]}] ~
  \item[\cs{db.get\_result\_by\_id(result\_id) $\to$ Optional[Tuple]}] ~
  \item[\cs{db.get\_result\_by\_alg(alg) $\to$ List[Tuple]}] ~
  \item[\cs{db.delete\_result(result\_alg) $\to$ None}] ~
  \item[\cs{db.delete\_all() $\to$ None}] ~
\end{description}

% =============================================================================
\section{Authors}
% =============================================================================

\begin{center}
\begin{tabular}{lll}
  \toprule
  Author & GitHub & Components \\
  \midrule
  Vladislav Gerda       & \href{https://github.com/hitfot}{@hitfot}       & OMP, CoSaMP, DCT, Metrics, Noise \\
  Grigory Demchenko     & \href{https://github.com/Pumukun}{@Pumukun}     & SP, BRGP, ImageCS \\
  Anastasia Kocherygina & \href{https://github.com/somniiium}{@somniiium} & --- \\
  Alina Sabitova        & \href{https://github.com/AlinaSAB}{@AlinaSAB}   & --- \\
  Mikhail Shibanov      & \href{https://github.com/Kar1ch}{@Kar1ch}       & Database (db.py) \\
  \bottomrule
\end{tabular}
\end{center}

\end{document}
